\section{Derivación bidireccional del exponente electrónico desde la ruta enriquecida}
\label{sec:cierre-electron-two-way-rev11}
\index{electrón!lector bidireccional}
\index{calendario de orden 108!exponente electrónico}

La ruta enriquecida del punto central publica dos representaciones
complementarias de su memoria: la palabra trítica completa
\(\gamma_{108}\in\mathbb F_3^{108}\) y el registro dodecafásico firmado
\(\tau_{12}\in\{-1,0,+1\}^{12}\). Los dos lectores definidos a continuación
se calculan exclusivamente desde esas representaciones, el calendario y las
constantes HMT ya construidas.

\subsection{Lector de desplazamientos sobre el calendario cíclico de orden 108}

Sea \(S\) el desplazamiento cíclico en \(\mathbb Z/108\mathbb Z\). Para
\(k\in\mathbb Z/108\mathbb Z\), definimos
\begin{equation}
 D_k(\Gamma):=
 \#\{t\in\mathbb Z/108\mathbb Z:\gamma_t\ne\gamma_{t+k}\}
 =\lVert(I-S^k)\gamma\rVert_0.
 \label{eq:lector-dk-ct108-rev11}
\end{equation}
El triple entero del lector es
\begin{equation}
 K_D(\Gamma)=
 \left(D_{27}+D_{36},\ D_1,\ \frac{D_4-D_3}{2}\right)
 =:(A_D,B_D,C_D),
 \label{eq:KD-rev11}
\end{equation}
y su evaluación escalar es
\begin{equation}
 \kappa_D(\Gamma)=A_D-\alpha_{\rm HMT}B_D+\Delta_4 C_D.
 \label{eq:kappaD-rev11}
\end{equation}
Los desplazamientos (27) y (36) son las elevaciones cronológicas de
\(90^\circ\) y \(120^\circ\) porque nueve pasos discretos representan
\(30^\circ\). El término \(D_1\) mide la frontera local inmediata y la
diferencia \(D_4-D_3\) compara las dos genealogías antes de su cociente de
altura. En el dominio de las 324 rutas se cumple
\(\gamma_{t+54}=\gamma_t\); todos los \(D_k\) son pares y \(C_D\) es entero.

\subsection{Lector dodecafásico de diferencias de correlación}

Para la palabra \(\tau=(\tau_r)_{r\in\mathbb Z/12\mathbb Z}\), sean
\[
 C_k^{\rm corr}(\tau)=\sum_{r=0}^{11}\tau_r\tau_{r+k},
 \qquad
 A_\ell(\tau)=\sum_{r=0}^{11}
 \left(\sum_{j=0}^{\ell-1}\tau_{r+j}\right)^2.
\]
La obstrucción genealógica primitiva y la memoria antipodal son
\begin{align}
 \Omega_{90\mid120}(\tau)
 &=4A_3-3A_4
 =-2C_1^{\rm corr}-4C_2^{\rm corr}-6C_3^{\rm corr},
 \label{eq:omega-two-way-rev11}\\
 P_6(\tau)&=\sum_{r=0}^{5}\tau_r\tau_{r+6}.
 \label{eq:p6-two-way-rev11}
\end{align}
El símbolo \(P_6\) evita la colisión con el ciclo corto del corredor de
constantes. El segundo triple y su evaluación son
\begin{equation}
 K_\Omega(\Gamma)=\bigl(\Omega_{90\mid120}(\tau_\Gamma),54,
 P_6(\tau_\Gamma)\bigr),
 \qquad
 \kappa_\Omega=\Omega_{90\mid120}-54\alpha_{\rm HMT}+P_6\Delta_4.
 \label{eq:kappa-omega-rev11}
\end{equation}

\HMTClaim{MD-R-ELECTRON-OMEGA-PRIMITIVE}
\begin{lemma}[Unicidad lineal integral primitiva]
\label{lem:unicidad-omega-rev11}
Sea \(L=uA_3+vA_4\), con \(u,v\in\mathbb Z\), un contraste que se anula
cuando las genealogías descienden a una misma altura,
\(A_3=3a\), \(A_4=4a\). Entonces
\((u,v)=n(4,-3)\). Fijado el signo \(90-120\),
\(\Omega_{90\mid120}=4A_3-3A_4\) es el generador primitivo.
\end{lemma}

\begin{proof}
La anulación para todo \(a\) exige \(3u+4v=0\). Como
\(\gcd(3,4)=1\), las soluciones enteras son
\((u,v)=n(4,-3)\). La primitividad corresponde a \(|n|=1\).
\end{proof}

El lema selecciona \(\Omega\) dentro de la clase de contrastes lineales
integrales de las energías \(A_3,A_4\). Su dominio de unicidad es esa clase
funcional; los invariantes no lineales pertenecen a un problema de selección
distinto.

\subsection{Evaluación de la ruta central}

La ruta electrónica \(\Gamma_e=(5,5,++)\) proporciona, con anterioridad al
contraste metrológico,
\[
 (D_1,D_3,D_4,D_{27},D_{36})=(54,56,68,48,32),
\]
y su registro dodecafásico firmado es
\[
 \tau_e=(+,+,+,-,-,-,+,+,+,-,-,-).
\]
Por tanto,
\[
 K_D(\Gamma_e)=(48+32,54,(68-56)/2)=(80,54,6),
\]
\[
 K_\Omega(\Gamma_e)=(80,54,6).
\]

\HMTClaim{MD-R-ELECTRON-TWOWAY-REDUCTION}
\begin{theorem}[Reducción electrónica común de los lectores bidireccionales]
\label{thm:reduccion-electron-two-way-rev11}
Los lectores \(K_D\) y \(K_\Omega\) coinciden sobre la ruta central y
determinan
\begin{equation}
 \kappa_e=80-54\alpha_{\rm HMT}+6\Delta_4.
 \label{eq:kappa-electron-demostrado-rev11}
\end{equation}
\end{theorem}

\begin{proof}
Las dos evaluaciones anteriores dan el mismo triple entero. La sustitución
de sus componentes en
\(A-\alpha_{\rm HMT}B+\Delta_4C\) prueba
\eqref{eq:kappa-electron-demostrado-rev11}. Toda cantidad empleada en esta
reducción pertenece a la ruta, a sus dos registros o a las constantes HMT
previamente construidas.
\end{proof}

\paragraph{Realización posterior en la carta de energía.}
Con \(R_{\rm act}=H_5/\etaRet\), la acción positiva sobre el basal
\(\mathfrak e_0\) produce
\begin{equation}
 \mathfrak e_e^\beta
 =\mathfrak e_0R_{\rm act}^{\kappa_e}.
 \label{eq:masa-electron-adimensional-two-way-rev11}
\end{equation}
Después de aplicar la carta vigente a MeV se obtiene
\begin{equation}
 \varsigma_{u_E}(\mathfrak e_e^\beta)\big|_{u_E=1\,\mathrm{MeV}}
 =0.5109989506900008086082571648\ldots\ \mathrm{MeV}.
 \label{eq:masa-electron-two-way-rev11}
\end{equation}
La primera igualdad pertenece al teorema interno condicionado por el basal y
la recta de acción; la segunda es su realización dimensional. El valor
externo se reserva para el contraste metrológico.

El entero (80) procede de \(D_{27}+D_{36}\) y, en el lector dodecafásico,
de \(\Omega_{90\mid120}\). La coincidencia con otros cardinales internos es
un control de compatibilidad y no sustituye esta procedencia operatoria. La
ruta de orientación \((--)\) es una traslación temporal de la ruta \((++)\)
y conserva el mismo perfil, conforme a la paridad de masa bajo conjugación.

\subsection{Separación de lectores y escalas de retorno}

La coincidencia numérica de componentes internos queda subordinada a su
tipo. En el punto central intervienen, entre otros, los objetos
\[
 R_{\rm atlas}(5,5)=(24,0,12,0,-12),
 \qquad
 L_{\kappa_0}(\Gamma_e^{++})=(24,0,0,0,-12),
\]
\[
 \begin{gathered}
 k_{\Delta_4}^{\rm dod},k_{\Delta_4}^{\rm pro}:
 \mathscr R_{324}\longrightarrow\mathbb Z,\\
 k_{\Delta_4}^{\rm dod}(\Gamma_e^{++})=12,
 \qquad
 k_{\Delta_4}^{\rm pro}(\Gamma_e^{++})=4,\\
 K_D(\Gamma_e^{++})=K_\Omega(\Gamma_e^{++})=(80,54,6).
 \end{gathered}
\]
Aquí \(\mathscr R_{324}\) denota el dominio de las 324 rutas orientadas.
El primero es una firma estática de familia; el segundo, un extractor de
genealogía; el tercer y el cuarto término son lectores torsionales definidos,
respectivamente, por la lectura dodecafásica y por la prolongación prospectiva;
el último objeto es el par de lectores bidireccionales. Ninguna igualdad
parcial identifica esos morfismos ni permite trasladar un coeficiente entre
lectores.

También se distinguen cuatro escalas temporales: \(27\) es el primer retorno
de clase, órbita y celda; \(54=D_1(\Gamma_e)\) es el retorno cinemático
observable dentro del ciclo; \(108\) cierra el calendario; y \(216\) cierra el estado cinemático
orientado en la extensión. Esta tipificación conserva separados retorno
posicional, observación y holonomía orientada.

\paragraph{Cadena causal del resultado electrónico}
La ruta y sus dos registros determinan \(K_D\), \(K_\Omega\), su reducción
común y la evaluación de \(\kappa_e\). La realización en MeV aplica después
la carta de energía; la masa externa interviene sólo en el contraste. El par
\(\mathbf K=(K_D,K_\Omega)\) es la salida operatoria completa de HMT sobre
cada ruta. En una fibra no central, su representación por un único escalar
físico exige un morfismo adicional entre el operador de fibra y la carta
dimensional; esta exigencia no modifica la ley interna ni autoriza a elegir
una celda por proximidad metrológica.
